\documentclass[12pt]{article}
\usepackage[utf8]{inputenc}
\usepackage[T1]{fontenc}
\usepackage{hyperref}
\title{Tchitchege Language Resources for K-12 Curriculum Development}
\author{Compiled by Claude Sonnet 5 (claude-sonnet-5) for Andr\'as Kornai}
\date{}
\begin{document}
\maketitle

\section{Introduction}
Tchitchege (Tsege; ISO 639-3 \texttt{tck}) is a Teke language of Gabon, also called Latsitsege or Lintsitsege. Wikipedia gives about 2,000 speakers (2003). The Endangered Languages entry places it near Mboua village on the road from Franceville to Boumango. It is a different language from Tsogo (\texttt{tsv}). This survey found very little material, and nothing designed for schools.

\section{Primary Educational Resources}
No primers, textbooks or school programmes turned up in the searches.

\section{Linguistic and Reference Materials}
\begin{itemize}
\item Wikipedia, Tsege language: \url{https://en.wikipedia.org/wiki/Tsege_language}
\item Endangered Languages Project entry: \url{https://www.endangeredlanguages.com/lang/5107}
\end{itemize}

\section{Religious and Community Resources}
\begin{itemize}
\item Joshua Project, Tchitchege in Gabon. The profile lists no Bible translation work, and its narrative text is not clearly sourced: \url{https://joshuaproject.net/people_groups/19532/GB}
\end{itemize}

\section{Digital and Audio Resources}
\begin{itemize}
\item Global Recordings Network, language entry for \texttt{tck}. I did not confirm what recordings it holds: \url{https://globalrecordings.net/en/language/tck}
\end{itemize}

\section{Curriculum Development Context}
Little published material exists. The language is thought to be spoken mainly by adults. Any curriculum would have to begin with documentation and community-led materials.
\end{document}